\documentclass[a4paper,12pt]{report}

\usepackage[
  % romanian / russian
  romanian,%
  % an: teza de an
  % licenta: teza de licenta
  % licenta_practica_1/2/3: practica la anul 1/2/3
  % master: teza de master
  % master_practica_2: practica de specialitate la anul 2 
  master%
]{config}

% design_jocurilor
% informatica
% informatica_aplicata
% informatica_stiinte_educatiei
\specialty{informatica_aplicata}

% Metadate.
\newcommand{\uniGroupName}{MIA2201}

\newcommand{\authorName}{Nume Prenume} % Fără patronimic
\newcommand{\conducatorName}{Nume Prenume}
\newcommand{\conducatorTitleRo}{asistent universitar} % doctor, conferențiar universitar

\newcommand{\docTitleRo}{Explorarea formatului PNG}
\newcommand{\docTitleEng}{Numele tezei}

\newcommand{\yearOfWriting}{\the\year}
\newcommand{\conferencesList}{Conferința Studențească, Editia NUMARUL-a, \yearOfWriting}
\newcommand{\github}{\url{https://github.com/USER/REPO}}
\newcommand{\outputDate}{\today}

% none - n-am folosit AI
% text - am folosit AI exclusiv pentru suport tehnic (corectare lingvistică, formatare, organizare preliminară)
% assisted - am folosit pentru altceva, explicat
\newcommand{\aiUsage}{text}
% Trebuie setat la o explicație, doar dacă ați ales "assisted".
% Dacă nu ați ales "assisted", aceasta va fi ignorată!
\newcommand{\aiAssistanceExplanation}{Am folosit pentru bla}

% Adăugați prescurtările aici.
% Adăugați doar prescurtările menționate.
% În viitor, mereu scrieți \gls{PNG} în loc de pur și simplu PNG.
\acro{PNG}{Portable Network Graphics}
\acro{PR}{Pull Request}
\acro[url_other_name]{URL}{Pull Request}
\acro{LONG TEST}{VERY LONG NAME VERY LONG NAME VERY LONG NAME VERY LONG NAME VERY LONG NAME VERY LONG NAME VERY LONG NAME VERY LONG NAME}

\begin{document}

\titlePage

\clearpage
\tableofcontents

% Doar pentru teză de master!
\annotationChapter{ro}

\textbf{la \docTypeRoNedefinit{} \enquote{\docTitleRo{}}, a studentului \authorName{}, grupa \uniGroupName{}, programul de studii \programulDeStudiiRo.}

\textbf{Structura tezei.}
Teza constă din: Introducere, \chapterCount{} capitole, Concluzii generale și recomandări, Bibliografie \bibliographyEntryCount{} titluri.
Textul de bază cuprinde \usefulPageCount{} de pagini și \anexeCount{} anexe.

\textbf{Cuvinte-cheie:}
% \acrshort înseamnă o referință la o prescurtare cu este.
% Face ca cuvântul să devină click-abil în PDF-ul generat.
\textit{\acrshort{PNG}, cuvânt 1, cuvânt 2}

\textbf{Actualitatea.}

Unde se folosesc tehnologiile, de ce e importantă temă.

\textbf{Scopul și obiectivele cercetării.}

Scopul = ce încercați să demonstrați prin această teză.

Obiectivele = lucrurile prin care demonstrați asta.

\textbf{Rezultatele preconizate și obținute.}
Acestea rezumă în:
(1) studierea ceva
(2) proiectarea ceva
(3) implementarea ceva.

\textbf{Problema importantă rezolvată.}
Problema importantă rezolvată în cercetare constă în ...

\textbf{Semnificația, noutatea și originalitatea.}
Lucrarea se deosebește prin ...

\textbf{Valoarea aplicativă, implementarea rezultatelor și publicarea.}
Ca rezultat s-a obținut ...

Toate codurile sursă a proiectului pot fi accesate pe GitHub după următorul link: \github.

Rezultatele obținute au fost raportate la \textbf{\conferencesList}\cite{self}.

% Doar pentru teză de master!
\annotationChapter{en}

\textbf{of the \docTypeEng{} \enquote{\docTitleEng{}} of the student \authorName{}, group \uniGroupName{}, \programulDeStudiiEng{} study program.}

\textbf{The structure of the thesis.}
The thesis is made up of: Introduction, \chapterCount{} chapters,
General conclusions and recommendations, Bibliography of \bibliographyEntryCount{} titles.
Base text takes up \usefulPageCount{} pages and has \anexeCount{} appendicies.

\textbf{Keywords:}
\textit{}

\textbf{Relevance.}

\textbf{Purpose and objectives of this research.}

\textbf{The planned and obtained results.}

\textbf{Important resolved problem.}

\textbf{Significance, novelty, and originality.}

\textbf{Applicative value, implementation of results, and publication.}

The entire source code can be obtained by accessing the GitHub repository
by the following \gls{url_other_name}: \github. 

The obtained results were reported at \textbf{\conferencesList}\cite{self}.

\acronymsChapter

\introChapter

\textbf{Actualitatea și importanța temei.}

Tot aceeași chestie ca și la adnotare, dar mai pe lung.

\textbf{Scopul și obiectivele.}

Tot aceeași.

\textbf{Criteriile de atingere a obiectivelor.}

Cum veți demonstra că fiecare obiectiv a fost realizat.

\textbf{Suportul metodologic și tehnologic.}

Ce fel de librării / instrumente / resurse ați folosit.

Ce fel de abordare ați luat.

Poate inspirațiile.

\textbf{Baza experimentală a cercetării.}

Ce date, scenarii, utilizatori, medii, seturi de testare sau alte condiții experimentale ați folosit.

\textbf{Noutatea științifică/originalitatea.}

Ce vă diferă teza / produs de alte aplicații / soluții / studii deja existente.

\textbf{Aportul concret al autorului.}

Care părți din rezultate sunt contribuția voastră personală.

\textbf{Semnificația teoretică a rezultatelor obținute.}

Ce clarifică, explică sau sistematizează rezultatele la nivel teoretic.

\textbf{Valoarea aplicativă.}

Pentru ce poate fi folosită aplicația.
Faptul că v-a ajutat să studiați aceste tehnologii nu ajunge, 
trebuie să mai specificați ceva în afară de asta.

\textbf{Aprobarea, publicarea și implementarea rezultatelor.}

Unde au fost raportate, publicate, testate sau implementate rezultatele.

Puteți face și astfel de listă:
\begin{itemize}
    \item Valoarea A
    \item Valoarea B
    \item Valoarea C
\end{itemize}

\textbf{Sumarul tezei.}

Primul capitol, \nameref{intro_chapter_title}, prezintă informații generale / teoretice despre ...

Al doilea capitol, \nameref{architecture_chapter_title}, schițează implementarea, ...

Al treilea capitol, \nameref{implementation_chapter_title}, urmează implementarea ...

\chapter{Denumirea capitolului teoretic}\label{intro_chapter_title}

\section{Sintaxa \LaTeX{}}

\subsection{Introducere}

Aici se demonstrează sintaxa \LaTeX{} de bază.

Evident că subcapitole de așa mărime nu sunt acceptabile, 
ar trebui să argumentați lucrurile mai mult.

De exemplu, s-ar putea scrie că sintaxa de bază este importantă
pentru a putea scrie chiar o lucrare primitivă în \LaTeX{}, și deoarece
scopul este să cunoască \LaTeX{} la nivel destul de avansat pentru a putea
scrie teza, este aceste informații sunt esențiale pentru a putea proceda la
niște idei mai adânce.


\subsection{Prescurtări}

\verb|\acronymsChapter| este comanda care inserează lista de prescurtări.
Când prescurtările au definiții lungi, acestea vor fi formatate automat frumos în tabel,
cu un număr minim de treceri la un rând nou.

Prescurtările în sine sunt definite înainte de începutul documentului, folosind \verb|\acro| sau \verb|\newacronym|.

\verb|\acro| este o comandă redefinită, folosită în șablon.
Primul parametru în paranteze pătrate permite acordarea unui identificator prescurtării, prin care
se va face referința la ea în fișierul sursă al lucrării.
Altfel, numele identificatorului va fi același ca prescurtarea.

\verb|\newacronym| necesită obligatoriu un identificator separat, ceea ce este rar necesar,
de aceea \verb|\acro| este mai ușor de utilizat.

Dacă prescurtările definite nu sunt menționate prin
\verb|\gls| sau \verb|\acrshort| sau un macro similar în lucrare,
ele nu vor apărea în lista de prescurtări.
Din acest motiv, \acrshort{LONG TEST} este menționat aici.

\subsection{Citarea surselor}

Citarea întregei surse se face după bucata de text care o menționează, așa\cite{gif_unusable_reason}.

Puteți face referință la pagini concrete, sau la paragrafe, așa\cite[12.2.]{png_spec}.

\subsection{Sintaxa listelor}

Neordonate. Folosiți pentru prezentarea unor informații legate:

\begin{itemize}
    \item
        Imagini cu paletă de culori de până la 256 de culori.
        
    \item
        Posibilitatea de a fi transmis în flux:
        fișierele pot fi citite și scrise în serie, permițând astfel utilizarea formatului
        PNG ca protocol de comunicare pentru generarea și afișarea imaginilor dinamic.

    \item
        Afișare progresivă: un fișier de imagine pregătit corespunzător poate fi
        afișat pe măsură ce este primit pe canalul de comunicare,
        furnizând foarte repede o imagine cu rezoluție redusă,
        urmată de îmbunătățirea graduală a detaliilor.

    \item
        Transparență: porțiuni ale imaginii pot fi marcate ca transparente,
        creând efectul unei imagini non-rectangulare.

    \item 
        Informații auxiliare: comentarii textuale și alte date pot fi
        stocate în interiorul fișierului de imagine.

    \item
        Complet independenți de hardware și platformă.

    \item
        Compresie eficientă, 100\% fără pierderi.
\end{itemize}

Ordonate. Folosiți de exemplu pentru o listă de pași consecutivi:

\begin{enumerate}
    \item 
        Imagini cu paletă de culori de până la 256 de culori.
    \item 
        Posibilitatea de a fi transmis în flux:
        fișierele pot fi citite și scrise în serie, permițând astfel utilizarea formatului
        PNG ca protocol de comunicare pentru generarea și afișarea imaginilor dinamic.
    \item 
        Afișare progresivă: un fișier de imagine pregătit corespunzător poate fi
        afișat pe măsură ce este primit pe canalul de comunicare,
        furnizând foarte repede o imagine cu rezoluție redusă,
        urmată de îmbunătățirea graduală a detaliilor.
    \item 
        Transparență: porțiuni ale imaginii pot fi marcate ca transparente,
        creând efectul unei imagini non-rectangulare.
\end{enumerate}


\subsection{Stilizarea}

\textbf{TEXT Bold Face = Bold}

\textit{TEXT ITalic = Italic}

\texttt{TEXT TeleType = font Monospace pentru cod}

Puteți să le \textbf{aplicați} \textit{într-o propoziție}, sau să le \texttt{\textit{\textbf{combinați}}}.

\subsection{Ghilimele}

Pentru citate scurte folosiți \verb|\enquote{...}|:
\enquote{acesta este un citat scurt}.

Puteți sublinia codul, de exemplu \verb|#include<iostrea> int a{~8};|.

Când textul literal trebuie folosit într-un argument fragil, de exemplu într-un titlu,
o inscripție sau un tabel, salvați-l cu \verb|\SaveVerb| și inserați-l cu \verb|\UseVerb|.

De asemenea, \verb|\verb| poate fi utilizat nu doar cu \verb!|! pentru a marca începutul și sfârșitul,
adică \verb!\verb|code|!, ci și, de exemplu, \verb|\verb!code!|.
Utilizați simbolul care nu este folosit în codul ce trebuie inserat ca atare.

\subsection{Tabele}

Când aveți un tabel, trebuie să dați referința la el.
Denumirea tabelului trebuie să se afle \textbf{deasupra} tabelului.
\Obiectul{pixel_order_table} prezintă un exemplu de inserare a tabelelor cu numerotare automată.

\insertTable[pixel_order_table]{Tabelul ordonării pixelilor\\a doua linie în titlu\\a treia linie}{%
  \begin{tabular}{c c c c c c c c}
      1 & 6 & 4 & 6 & 2 & 6 & 4 & 6 \\
      7 & 7 & 7 & 7 & 7 & 7 & 7 & 7 \\
      5 & 6 & 5 & 6 & 5 & 6 & 5 & 6 \\
      7 & 7 & 7 & 7 & 7 & 7 & 7 & 7 \\
      3 & 6 & 4 & 6 & 3 & 6 & 4 & 6 \\
      7 & 7 & 7 & 7 & 7 & 7 & 7 & 7 \\
      5 & 6 & 5 & 6 & 5 & 6 & 5 & 6 \\
      7 & 7 & 7 & 7 & 7 & 7 & 7 & 7 \\
  \end{tabular}
}


\SaveVerb{mintedLibraryName}|minted|

\section[Librăria minted pentru cod]{Librăria \UseVerb{mintedLibraryName} pentru cod}

\subsection{Informații generale}

Librăria \verb|minted| se folosește pentru a stiliza codul.
Blocurile de cod se inserează prin helper-ele template-ului, la fel ca imaginile și tabelele,
pentru a primi numerotare, inscripție și plasare automată.

\subsection{Inserarea directă a unui bloc de cod}

\Obiectul{code_direct_example} prezintă inserarea directă a unui bloc de cod.

\begin{code}[code_direct_example]{cpp}{Exemplu de bloc de cod inserat direct}
// orice cod aici
int main()
{
    std::cout << "hello";
}
\end{code}

\subsection{Inserarea unui fișier întreg}

\Obiectul{../src/sourcefile.zig} inserează un fișier întreg.

\insertCodeFile{zig}{../src/sourcefile.zig}{Fișier sursă complet}

\subsection{Inserarea unei bucăți din fișier}

\Obiectul{code_range_example} inserează numai intervalul de linii transmis prin opțiunile \verb|minted|.

\insertCodeFile[code_range_example][firstline=2,lastline=5]{zig}{../src/sourcefile.zig}{Interval de linii din fișier}

\subsection{Inserarea unui segment din fișier}

Această funcționalitate este implementată cu ajutorul scriptului \texttt{findSegment.py}.

\verb|minted| oferă o funcționalitate similară (\verb|rangestartafterstring|, \verb|rangestopbeforestring|),
însă nu suportă afișarea numerelor de linie din fișier și inserează o linie goală la începutul segmentului.

Limbajul de programare este determinat pe baza extensiei fișierului.
Pentru a adăuga extensii care nu sunt încă suportate, modificați \texttt{findSegment.py}.

\Obiectul{code_segment_example} inserează un segment delimitat prin comentarii în fișierul sursă.

\insertCodeSegment[code_segment_example]{../src/sourcefile.zig}{example}{Segment marcat din fișier}

\Obiectul{code_segment_autogobble_example} inserează un segment indentat suplimentar.
Opțiunea \verb|autogobble| elimină indentarea comună din fața codului afișat.

\insertCodeSegment[code_segment_autogobble_example]{../src/sourcefile.zig}{autogobble_example}{Segment cu indentare eliminată automat}

\subsection{Inserarea textului pe mai multe rânduri ca cod}

\Obiectul{code_text_example} inserează text simplu ca bloc de cod.

\begin{code}[code_text_example]{text}{Text pe mai multe rânduri inserat ca cod}
Regular text
Multiline
With any characters |||
\end{code}

Pentru a insera doar text, fără numărul de rând, este necesar să adăugați \verb|[linenos=false]|

\Obiectul{code_text_no_lines_example} ascunde numerele de rând prin opțiunile \verb|minted|.

\begin{code}[code_text_no_lines_example][linenos=false]{text}{Text fără numere de rând}
Regular text
Multiline
With any characters |||
\end{code}

\subsection{Inserarea codului literal în tabele}

Pentru fragmente scurte de cod care trebuie să apară într-un tabel sau în inscripția lui,
salvați textul literal înainte de tabel și inserați-l cu \verb|\UseVerb|.

\SaveVerb{tableHttpRoute}|GET /users/{id}?q=a&b%20c|
\SaveVerb{tableJsonConfig}|{"path": "C:\tmp\{id}", "enabled": true}|

\Obiectul{inline_code_table_example} prezintă inserarea unor fragmente salvate cu \verb|\SaveVerb|.

\insertTable[inline_code_table_example]{Fragmente de cod salvate: \UseVerb{tableHttpRoute}}{%
  \begin{tabularx}{\textwidth}{l X}
    \toprule
    Context & Fragment \\
    \midrule
    Rută HTTP & \UseVerb{tableHttpRoute} \\
    Configurare JSON & \UseVerb{tableJsonConfig} \\
    \bottomrule
  \end{tabularx}
}

\section{Imagini}

\subsection{Introducere}

Acest șablon propune doar o funcție pentru imagini, 
care înregistrează automat imaginea ca figură.

Imaginea va fi plasată \textit{nu imediat} după mențiune,
ci va apărea automat în locul cel mai apropiat de punctul de mențiune,
astfel încât să minimizeze spațiul gol
lăsat atunci când imaginea nu încape.

Puteți folosi altceva dacă nu vă aranjează.
Pentru standardizare, însă, ar fi bine să vă propuneți metoda cu un \gls{PR} ca s-o adăugați în șablon.

\subsection{O imagine centrată}

Când dați o imagine, la ea numaidecât să faceți și referința.

În \obiectul{interface_sketch.png} este prezentată o schiță a interfeței grafice.

\insertImage{interface_sketch.png}{Schița interfeței grafice}

Această metodă limitează lățimea și înălțimea maximă ale imaginilor,
fără a le permite să iasă în afara paginii și fără a le modifica raportul de aspect.
De exemplu, în \obiectul{wide.png} o imagine lată nu va depăși limitele paginii,
iar în \obiectul{long.png} nu va ieși în afara paginii.

\insertImage{wide.png}{Bandă lată}

\insertImage{long.png}{Bandă lungă}

\section{Anexe}

\subsection{Alegerea între anexă și cod direct}

Atunci când conținutul codului nu este strict esențial 
pentru a-l explica la nivel înalt în text.

Încă atunci, când codul este prea mare. 
Încercați să includeți direct în text cam 50 de linii maxim.

\subsection{Cum se face referința}

A se vedea anexa \ref{appendix:example_min}.

În anexa \ref{unnumbered_appendix} este prezentată o anexă fără numerotarea elementelor din interior.

În anexa \ref{numbered_appendix} este prezentată o anexă cu numerotarea elementelor din interior,
și anume:

\Obiectele{cat_1}{cat_2} demonstrează pisicile.
\begin{itemize}
  \item \Obiectul{cat_1} demonstrează pisica numărul 1
  \item \Obiectul{cat_2} demonstrează pisica numărul 2
\end{itemize}

În anexa \ref{whole_appendix_reference} este prezentat un exemplu în care
referința la întreaga anexă acoperă imaginile, tabelele și blocurile de cod numerotate din ea.

\chapterConclusionSection{intro_chapter_title}

Ce s-a discutat în acest capitol, pe scurt.


\chapter{Proiectarea aplicației}\label{architecture_chapter_title}

\section{Introducere}

Aici includeți informații despre structura aplicației la nivel înalt, 
poate niște diagrame, de ce ați folosit tehnologiile pe care le-ați folosit, etc.

\chapterConclusionSection{architecture_chapter_title}

\chapter{Implementarea Sistemului}\label{implementation_chapter_title}

\section{Introducere}

Aici includeți mai multe detalii despre module concrete din sistemul implementat.
Ce probleme ați întâlnit și cum le-ați depășit.

Cum ați folosit funcționalitățile importante din librării.

Prezentați aplicația cu imagini, ori direct în text, 
dacă ele sunt importante pentru înțelegere și mai specifice,
ori cu referințe la anexe.

\chapterConclusionSection{implementation_chapter_title}

În acest capitol s-au discutat ...

\finalConclusionChapter

Ce s-a discutat pe scurt, reamintiți cele mai interesante momente.
Mai adăugați ce ar putea fi îmbunătățit sau unde poate fi folosit produsul în viitor.

Un exemplu al unei teze complet:
\begin{itemize}
  \item Sursa (programul, textul tezei în \LaTeX{}): \url{https://github.com/AntonC9018/thesis-png}
  \item PDF-ul: \url{https://drive.google.com/file/d/1ZiGQt6PvUm3FGY8oyIlVsERc4PvZHJu0/view?usp=drive_link}
\end{itemize}

Acestea 2 de mai jos le includeți mereu:

Toate codurile sursă, inclusiv codurile de program și textul lucrării
în forma înainte de randare, pot fi accesate pe GitHub după următorul link: \github.

Rezultatele obținute au fost raportate la \textbf{\conferencesList}\cite{self}.

\bibliographyChapter

\appendixChapter

\section{Funcția min într-un program test JavaScript}\label{appendix:example_min}
\insertCodeFile*{js}{../src/appendix_example.js}

\section{Fără numerotare}\label{unnumbered_appendix}

\insertImage*{cat_1.jpg}{Pisică}

Când referința directă este doar la întreaga anexă,
nu la imaginile, tabelele sau blocurile de cod individuale din interior,
acestea nu trebuie neapărat numerotate.

\insertTable*{Exemplu de tabel fără numerotare}{%
  \begin{tabular}{c c c}
    Hello & 1 & 2 \\
    A & B & C \\
    D & E & F
  \end{tabular}
}

Pentru cod direct nenumerotat în anexe se poate folosi mediul \verb|codeWithoutLabel|.

\begin{codeWithoutLabel}{js}
function min(a, b) {
  return a < b ? a : b;
}
\end{codeWithoutLabel}

\section{Imagine cu numerotare}\label{numbered_appendix}

Are sens să procedați astfel dacă în anexă există mai mult de un element.

\insertImage[cat_1]{cat_1.jpg}{Pisica 1}

\insertImage[cat_2]{cat_2.jpg}{Pisica 2}

\section{Referință la întreaga anexă}\label{whole_appendix_reference}

În acest caz, elementele numerotate din anexă sunt acoperite de referința la anexă.

\insertImage[whole_appendix_cat]{cat_1.jpg}{Pisică acoperită de referința la anexă}

\insertTable[whole_appendix_table]{Tabel acoperit de referința la anexă}{%
  \begin{tabular}{c c}
    Element & Valoare \\
    Imagine & 1 \\
    Tabel & 2
  \end{tabular}
}

\insertCodeFile[whole_appendix_code][]{js}{../src/appendix_example.js}{Cod acoperit de referința la anexă}

\ifdefstring{\documentIsPractice}{false}{\declarationPage{}}{}

\end{document}
% vim: fdm=syntax
